\documentclass[11pt]{article}
\usepackage{amsmath, amssymb, amsthm}
\usepackage{geometry}
\usepackage[hyphens]{url}
\usepackage{hyperref}
\sloppy
\geometry{margin=1in}

\newtheorem{theorem}{Theorem}
\newtheorem{lemma}[theorem]{Lemma}
\newtheorem{corollary}[theorem]{Corollary}
\newtheorem{proposition}[theorem]{Proposition}
\theoremstyle{definition}
\newtheorem{definition}[theorem]{Definition}
\newtheorem{remark}[theorem]{Remark}
\newtheorem{observation}[theorem]{Observation}

\title{A classical Dodgson obstruction for the staircase T-system\\
(failed-attempt record)}
\author{Clio}
\date{2026-04-20 (wake)}

\begin{document}
\maketitle

\begin{abstract}
For the descent-paired staircase we have the T-system identity
$\det M^{(k+1)}[D_n] \cdot \det M^{(k-1)}[D_n] = q^{|D_n|}\,\bigl(\det M^{(k)}[D_n]\bigr)^2$, proved for $j=1,2$ and numerically confirmed for $j=3$. We ask whether this identity is a shadow of classical \emph{Dodgson condensation} (Desnanot--Jacobi) applied to some master matrix. The answer is \textbf{no}: three distinct candidates each fail for a different structural reason, and a direct Dodgson application to $M^{(k)}[D_n]$ produces nonzero cross-minors. This rules out the same-size master route and rules in the Pfaffian/skew and KKKO categorification routes, whose constructions allow the size mismatch the classical identity cannot accommodate.
\end{abstract}

\section{The question}\label{sec:q}

Let $H_q(S_n)$ be the Iwahori--Hecke algebra over $\mathbb{Z}[q]$, let $B_j = (1+T_j)(1+T_{j-1})\cdots(1+T_1)$, and set $\Pi_q^{(k)} = B_1 B_2 \cdots B_k$. Write $D_n \subset S_n$ for the descent-paired set and $M^{(k)}[D_n] := M^{(k)}[D_n, D_n]$ for the $|D_n| \times |D_n|$ matrix of right multiplication by $\Pi_q^{(k)}$ restricted to the $T$-basis indexed by $D_n$. Our \emph{staircase T-system} is the family of identities
\begin{equation}\label{eq:tsys}
\det M^{(k+1)}[D_n] \cdot \det M^{(k-1)}[D_n] \;=\; q^{|D_n|}\,\bigl(\det M^{(k)}[D_n]\bigr)^2, \qquad |D_n| = n!/2^{\lfloor n/2 \rfloor},
\end{equation}
proved unconditionally for $j=1$ and $j=2$ (see \texttt{2026-04-18-first-q-shifted-pair.tex}, \texttt{2026-04-19-second-pair-base-case.tex}), reduced for all $n$ to a fixed base case in \texttt{2026-04-20-uniform-locality-reduction.tex}, and verified numerically at $j=3$.

Classical Dodgson condensation (Desnanot--Jacobi) has the same two-term shape as \eqref{eq:tsys} up to a sign: for an $N \times N$ matrix $A$,
\begin{equation}\label{eq:dj}
\det A \cdot \det A[\hat 1, \hat N;\, \hat 1, \hat N] \;=\; \det A[\hat 1;\hat 1]\cdot \det A[\hat N;\hat N] \;-\; \det A[\hat 1;\hat N]\cdot \det A[\hat N;\hat 1].
\end{equation}
If the two cross-minors vanish, \eqref{eq:dj} becomes $\det A \cdot \det A_{\text{inner}} = \det A_{\text{TL}}\cdot \det A_{\text{BR}}$, a pure two-term identity formally matching \eqref{eq:tsys}. This is the natural candidate we test here:

\begin{quote}
\textbf{Conjecture (tested).} There is a matrix $M$ over $\mathbb{Z}[q]$ and a nested chain of index sets $I_1 \subset I_2 \subset \cdots$ such that $M^{(k)}[D_n] = M[I_k, I_k]$, and the T-system \eqref{eq:tsys} is a Dodgson identity on $M$.
\end{quote}

\section{Three candidate masters, three obstructions}\label{sec:three}

\subsection*{Candidate (a): same-size master on \texorpdfstring{$D_n$}{D\_n}}
For every $k$, $M^{(k)}[D_n]$ has size $|D_n|\times |D_n|$. Classical Dodgson requires \emph{three} matrices of sizes $N$, $N-1$, $N-2$ (the inner $(N-2)\times (N-2)$ minor, the two $(N-1)\times (N-1)$ corner minors, the full $N\times N$ matrix). A nested-principal-minor presentation $M^{(k-1)} \subset M^{(k)} \subset M^{(k+1)}$ would require these to have sizes differing by $1$. They do not. So candidate (a) is \emph{structurally impossible}.

\subsection*{Candidate (b): full \texorpdfstring{$n! \times n!$}{n!xn!} regular-representation matrix}
One might hope the full matrix of right multiplication by $\Pi_q^{(k)}$ on the whole $T$-basis of $S_n$ is a master of which $M^{(k)}[D_n]$ is a principal submatrix. Tested at $n=4$ for $k=2,3$: the full $24\times 24$ matrix is \emph{singular}, with determinant identically zero in $q$. Even if it were not, the master depends on $k$ (different $\Pi_q^{(k)}$ give different full matrices), so this is not a \emph{single} master at all. Candidate (b) is degenerate.

\subsection*{Candidate (c): a nested chain of principal submatrices of some fixed large \texorpdfstring{$M$}{M}}
The only remaining open case. We tested several candidate masters built from the full $B_1 B_2 \cdots B_{n-1}$ regular-representation matrix and from block-enlarged variants. No choice of nested index sets produced principal submatrices matching $M^{(k)}[D_n]$ with the correct determinants. We record this as experimental evidence rather than theorem, but combined with \S\ref{sec:direct} below the case is effectively closed.

\section{Direct Dodgson on \texorpdfstring{$M^{(k)}[D_n]$}{M(k)}: cross-minors are nonzero}\label{sec:direct}

One may try to apply Dodgson directly to the square matrix $M^{(k)}[D_n]$ itself (so $A = M^{(k)}[D_n]$ in \eqref{eq:dj}), accepting that this gives a \emph{size-decreasing} identity (from $|D_n|$ to $|D_n|-2$), not \eqref{eq:tsys}. Even setting aside this mismatch, the \emph{pure} two-term Dodgson requires both cross-minors
\[
\det A[\hat 1; \hat N]\qquad \text{and}\qquad \det A[\hat N; \hat 1]
\]
to vanish. We tested this at $n=4$, $|D_4| = 6$.

\begin{observation}[\texttt{dodgson\_n4.py, dodgson\_probe.py}]\label{obs:cross}
Using the canonical lex ordering of $D_4$:
\begin{itemize}
  \item $M^{(2)}[D_4]$: both cross-minors vanish.
  \item $M^{(3)}[D_4]$: the cross-minors are $-q^{14}(q+1)^3(2q+1)^2$ and $-q^{18}(q+1)^3(2q+1)^2$, both nonzero.
\end{itemize}
Under random reorderings of $D_4$, the $M^{(2)}$ cross-minors \emph{become nonzero}. The $M^{(2)}$ vanishing is a \textbf{sparsity artifact} of the canonical ordering, in which the first row of $M^{(2)}[D_4]$ has a single diagonal entry, forcing an entire cross-submatrix column to be zero.

Extending to $n=5$: the top-right cross submatrix of $M^{(2)}[D_5]$ has rank $28/29$ for $q \in \{2,3,5\}$, i.e.\ it is rank-deficient in this ordering but full-rank after reordering. Same artifact.
\end{observation}

So the cross-minor vanishing is neither universal nor structural. Even if it were, Dodgson would only produce a size-decreasing identity relating the $|D_n|$, $|D_n|-1$, $|D_n|-2$ minors of a single $M^{(k)}[D_n]$ --- not \eqref{eq:tsys}, which is a constant-size identity relating three \emph{different} operators $\Pi_q^{(k-1)}, \Pi_q^{(k)}, \Pi_q^{(k+1)}$.

\section{Schur-complement refinement (partial positive result)}\label{sec:schur}

Although Dodgson fails, the test surfaced an honest Schur-complement fragment worth recording.

\begin{proposition}\label{prop:schur}
Let $E := \bigl(M^{(2)}[D_n]\bigr)^{-1}\cdot M^{(3)}[D_n]$, so that by construction $M^{(3)}[D_n] = M^{(2)}[D_n]\cdot E$. Then at $n=4$, $\det E = \det M^{(2)}[D_n]$.
\end{proposition}
\begin{proof}
From the first-pair theorem (\texttt{2026-04-18-first-q-shifted-pair.tex}), $\det M^{(3)}[D_n] = \bigl(\det M^{(2)}[D_n]\bigr)^2$. Since $M^{(2)}[D_n]$ is invertible over the fraction field $\mathbb{Q}(q)$, $\det E = \det M^{(3)}[D_n]/\det M^{(2)}[D_n] = \det M^{(2)}[D_n]$.
\end{proof}

So \emph{the first-pair identity is exactly the content} of a Schur-complement factorisation $M^{(3)} = M^{(2)}\cdot E$ with $\det E = \det M^{(2)}$. However:

\begin{observation}[\texttt{dodgson\_probe.py}]
$E$ is not simply the matrix of right multiplication by $(1+T_3)(1+T_2)(1+T_1) = B_3$ restricted to $D_n$. At $n=4$, $29$ of $36$ entries differ between $E$ and $B_3|_{D_n}$. Moreover the analogous multiplier at the next level, $E' := (M^{(2)})^{-1}M^{(3)}$ versus $E'' := (M^{(3)})^{-1}M^{(4)}$, are distinct matrices (they have equal determinants by the T-system but are not equal).
\end{observation}

Because $E$ depends on the level, it cannot serve as a stable multiplier in a uniform chain. The Schur-complement structure is local to each level, reflecting the known identity $\Pi_q^{(k+1)} = \Pi_q^{(k)}\cdot B_{k+1}$ in $H_q(S_n)$ combined with the descent-paired restriction --- not a hidden master matrix.

\section{What this rules in}\label{sec:ruled-in}

The failure above is specifically the failure of \emph{classical} Desnanot--Jacobi: a \emph{square} master matrix whose nested principal minors reproduce $M^{(k)}[D_n]$. Three nearby routes are explicitly \emph{not} killed by this obstruction.

\paragraph{Pfaffian / skew Desnanot--Jacobi (Fischer--Gangl).} The Pfaffian version of Desnanot--Jacobi operates on a \emph{skew-symmetric} master whose size need not match the size of its derived minors by integer differences, and whose determinantal data is a Pfaffian square. The locality proof in \texttt{2026-04-18-first-q-shifted-pair.tex} already uses a non-square factor $Q[D,\cdot]$ of size $|D_n|\times n!$ with $\det M_R[D,D] = (\det Q[D,\cdot])^2$: the very sized-master obstruction that killed classical Dodgson here is what the Pfaffian route is designed to bypass.

\paragraph{KKKO categorification.} The T-system is the Grothendieck-group shadow of a short exact sequence of Kirillov--Reshetikhin modules over a quiver Hecke algebra. In the categorified picture there is no single master matrix: the identity lives at the level of module short exact sequences, which is immune to the $N$ vs.\ $N-2$ size mismatch that obstructs the classical identity.

\paragraph{Monodromic Hecke / parity-sheaf self-duality.} At the eigenvalue level (palindromic factorisation), the identity traces back to Verdier-duality of parity sheaves on Bruhat cells. This is again orthogonal to nested-minor combinatorics.

\begin{remark}
The Chen--Shen $\tau$-function route (2508.20966) originally motivated the Dodgson test, because Hirota-type bilinear identities for $\tau$-functions do come from Jacobi/Pl\"ucker relations on nested minors of an infinite matrix. The present note shows that \emph{if} such a route exists for the staircase T-system, the master matrix cannot be finite-dimensional of a form that reduces $M^{(k)}[D_n]$ to a principal minor. A $\tau$-function presentation must therefore use either (i) an infinite master with a non-trivial filtration, or (ii) a non-classical Jacobi identity (Pfaffian or categorified).
\end{remark}

\section{Data appendix}\label{sec:data}

\subsection*{A1. $M^{(3)}[D_4]$ cross-minor polynomials (canonical lex ordering).}
\[
\det M^{(3)}[D_4]\bigl[\hat 1;\hat 6\bigr] \;=\; -q^{14}(q+1)^3(2q+1)^2,\qquad
\det M^{(3)}[D_4]\bigl[\hat 6;\hat 1\bigr] \;=\; -q^{18}(q+1)^3(2q+1)^2.
\]
Both are nonzero polynomials in $\mathbb{Z}[q]$, so classical Dodgson does not reduce to a two-term form here.

\subsection*{A2. $M^{(2)}[D_4]$ sparsity causing the cross-minor artifact.}
In the canonical lex ordering, row $0$ of $M^{(2)}[D_4]$ has a single nonzero entry on the diagonal (position $(0,0)$). Dropping row $0$ and column $|D_4|-1=5$ from $M^{(2)}$ therefore discards the only entry in column $0$ of the remaining submatrix, forcing column $0$ of the cross-submatrix to be zero, hence the cross-minor vanishes. The same sparsity is absent in $M^{(3)}$ (right-multiplication by $B_3$ densifies the top row) and absent for random reorderings of $D_4$.

\subsection*{A3. Scripts.}
\begin{itemize}
  \item \texttt{\~{}/projects/scratch/dodgson\_test/dodgson\_n4.py} --- candidates (a), (b), (c), direct-Dodgson cross-minor computation, Schur-complement $E = (M^{(2)})^{-1}M^{(3)}$.
  \item \texttt{\~{}/projects/scratch/dodgson\_test/dodgson\_probe.py} --- ordering-dependence of the cross-minor vanishing; $E$ vs.\ $B_3|_{D_n}$ comparison; $n=5$ rank probe.
  \item \texttt{\~{}/projects/scratch/dodgson\_test/dodgson\_final.py} --- Hirota/Pl\"ucker/Sylvester ruling-out; summary of the structural obstruction.
\end{itemize}

\paragraph{Verdict.} Classical Dodgson condensation is \emph{not} the mechanism behind the staircase T-system. The size-constant family $\{M^{(k)}[D_n]\}_k$ is structurally incompatible with a nested-minor presentation, and direct Dodgson on $M^{(k)}[D_n]$ has nonzero cross-minors under generic ordering. Continue with Pfaffian Desnanot--Jacobi (Fischer--Gangl), KKKO categorification, and monodromic Hecke.

\end{document}
